\section{Requisiti di sistema}
\subsection{Browser Supportati}

Di seguito viene fornito un breve elenco delle versioni minime dei browser sui quali il funzionamento del nostro plugin è garantito:
\begin{itemize}
	
	\item Google Chrome v.55
	\item Mozilla Firefox v.50
	\item Safari v.10
	\item Internet Explorer v.11
	
\end{itemize}

Per un funzionamento corretto del plugin è necessario che sia abilitato Javascript.


\section{Manuale Utente}
\subsection{Mappa Topologica}
\label{sec:mtopologica}
\subsubsection{Interfaccia}
Questo plugin si presenta come una mappa nella quale vengono visualizzati i nodi di cui è composto il sistema. 
I nodi vengono disegnati graficamente in modo diverso a seconda del tipo di componente che rappresentano, ad esempio i server sono rappresentati con un'icona circolare di colore rosso, mentre i database sono rappresentati da un cilindro di colore blu e nero. Ogni nodo è affiancato dal nome del componente. 
\begin{figure}[H]
	\centering 
	\includegraphics[width=0.08\textwidth]{Images/server.png}
	\caption{icona utilizzata per rappresentare i server}
	\label{server} 
\end{figure}
\begin{figure}[H]
	\centering 
	\includegraphics[width=0.08\textwidth]{Images/database.png}
	\caption{icona utilizzata per rappresentare i database}
	\label{database} 
\end{figure}

\subsubsection{Funzionalità}
    \paragraph {Drag and Drop} \Spazio
Questa funzionalità permette di spostare i nodi della mappa nella posizione desiderata purchè sia all'interno dell'area del plugin il cui perimetro è evidenziato da un rettangolo di colore grigio.
Per spostare un nodo è necessario: 
\begin{enumerate}
	
	\item Portare il cursore del mouse sopra al nodo che si intende muovere;
	\item Cliccare e tenere premuto sul nodo;
	\item Spostare il nodo nella posizione desiderata.
	
\end{enumerate}

\paragraph{Zoom} \Spazio
Questà funzionalità permette all'utente di effettuare le operazioni di ingrandimento e rimpicciolimento delle dimensioni di visualizzazione della mappa. In particolare, è possibile ingrandire una specifica area della mappa.
Per effettuare questa operazione è necessario:

\begin{enumerate}
	
	\item Portare il cursore del mouse nel punto della mappa in cui si vuole effettuare l'operazione di zoom;
	\item Muovere la rotellina del mouse in avanti se si intende ingrandire, all'indietro se si intede rimpicciolire. Nel caso non si disponga della rotellina sul mouse, utilizzare le funzionalità presenti nella tastiera per effettuare lo zoom.
	
\end{enumerate}

\paragraph{Tempo di risposta} \Spazio
Questa funzionalità permette di visualizzare il tempo medio di risposta impiegato per effettuare richieste tra due nodi del sistema.
E' possibile visualizzarlo in corrispondenza della linea che collega due componenti.
Se il numero è colorato in rosso, significa che il tempo medio di risposta è superiore rispetto ad una soglia standard definita.




\subsection{Stack Trace}
\label{sec:strace}
\subsubsection{Interfaccia}
Questo plugin è rappresentanto mediante una lista nella quale sono presenti tutte le richieste effettuate dal sistema. Per ogni richiesta si possono visualizzare diverse informazioni, nello specifico:
   
    \begin{itemize}
    	
    	\item \textbf{Numero della richiesta:} Numero progressivo associato alla richiesta;
    	\item \textbf{Richiesta:} Nome della richiesta;
    	\item \textbf{Execution Time:} Tempo totale di esecuzione della richiesta;
    	\item \textbf{Timestamp:} Indica l'istante temporale in cui è stata effettuata la richiesta;
    	\item \textbf{Errore:} Indica se durante l'esecuzione della richiesta si è presentato qualche errore.
    	
    	\end{itemize}

\subsubsection{Funzionalità}
\paragraph {Dettaglio richiesta} \Spazio
Questa funzionalità permette di raggiungere un maggiore grado di dettaglio rispetto alle informazioni inizialmente rappresentate sulla tabella nella riga corrispondente alla richiesta di nostro interesse.
Per visualizzare maggiori informazioni riguardo una richiesta, l'utente deve cliccare sulla riga corrispondente.
Nel popup che appare è possibile, mediante il pulsante "switch", navigare nella \emph{Call Tree} o nella \emph{Query List} a seconda delle necessità.

\subparagraph {Call Tree} \Spazio
Rappresenta tutte le funzioni chiamate dal sistema per effettuare la richiesta dell'utente.
Ogni funzione è stata disegnata in modo identanto così da riuscire a visualizzare la gerarchia di annidamento per ogni metodo padre.
E' possibile espandere/restringere ogni gerarchia di funzioni cliccando su di essa, in questo modo, tutta la gerarchia annidata verrà visualizzata/nascosta.

In corrispondenza di ogni funzione è possibile visualizzare il tempo di esecuzione.

   

\subparagraph {Query List} \Spazio
Rappresenta tutte le query effettuate nel database dal sistema per fornire una risposta alla richiesta dell'utente.
    \begin{itemize}
	
	\item \textbf{Numero della query:} Numero progressivo associato alla query;
	\item \textbf{Query:} Query effettuata;
	\item \textbf{Database:} Nome del database nel quale è stata effetuata la query;
	\item \textbf{Execution Time:} Tempo impiegato dalla query per fornire un risultato;
	\item \textbf{Timestamp:} Indica l'instante temporale in cui è stata effettuata la query.
	
\end{itemize}







